% ═══════════════════════════════════════════════════════════════
%  PART VI — APRIL 8, 2026: CASCADE BREAKTHROUGHS
% ═══════════════════════════════════════════════════════════════
\clearpage
\part{The Cascade: From Threshold Corrections to the Meaning of $\alpha$}

\section{Threshold Corrections Fix the RG Running}
\label{sec:threshold}

The na\"ive prediction $\alpha_1^{-1}=\alpha_2^{-1}=\alpha_3^{-1}=f=24$
at the Planck scale fails for the pure Standard Model.
The resolution: the W(3,3) finite spectral triple contributes
\textbf{threshold corrections} from its three distinguished substructures.

\begin{theorem}[Corrected Gauge Coupling Boundary Conditions]
\label{thm:threshold}
\begin{equation}
\alpha_i^{-1}(M_{\mathrm{Pl}}) = f + \Delta_i, \qquad
f = 24,
\end{equation}
where the threshold corrections are:
\begin{equation}
\boxed{
\Delta_1 = \Theta = 10, \quad
\Delta_2 = f = 24, \quad
\Delta_3 = q^3 = 27
}
\end{equation}
giving $\alpha_i^{-1}(M_{\mathrm{Pl}}) = (34,\;48,\;51)$.
\end{theorem}

\noindent\textbf{Comparison with observation} (running observed $M_Z$ values
upward with 1-loop SM $\beta$-functions):
\begin{center}
\begin{tabular}{@{}lccc@{}}
\toprule
& W(3,3) prediction & Upward from data & Error \\
\midrule
$\alpha_1^{-1}(M_{\mathrm{Pl}})$ & 34 & 34.3 & 0.9\% \\
$\alpha_2^{-1}(M_{\mathrm{Pl}})$ & 48 & 48.6 & 1.3\% \\
$\alpha_3^{-1}(M_{\mathrm{Pl}})$ & 51 & 50.6 & 0.7\% \\
\bottomrule
\end{tabular}
\end{center}

Running back down to $M_Z$: $\sin^2\theta_W = 0.228$
(observed $0.2312$, error 1.2\%).

\begin{theorem}[The $E_7$ Sum Rule]
\begin{equation}
\boxed{\sum_{i=1}^{3}\alpha_i^{-1}(M_{\mathrm{Pl}}) = 34+48+51 = 133 = \dim(E_7)}
\end{equation}
This arises from the maximal subgroup decomposition
$E_8\supset E_7\times\SU(2)$:
\begin{equation}
248 = (133,\mathbf{1})\oplus(\mathbf{1},3)\oplus(56,\mathbf{2}).
\end{equation}
\end{theorem}


\section{Electroweak Symmetry Breaking as a $G_2$ Shift}
\label{sec:ewsb-g2}

\begin{theorem}[Pre-EWSB Electroweak Unification]
Before electroweak symmetry breaking:
\begin{equation}
\alpha_1^{-1} = \alpha_2^{0\,-1} = \lambda(\Phithree+\mu) = 34,
\qquad
\sin^2\theta_W^{0} = \frac{3}{8}.
\end{equation}
The standard $\SU(5)$ GUT prediction $\sin^2\theta_W = 3/8$
holds \emph{before} the Higgs mechanism acts.
\end{theorem}

\begin{theorem}[EWSB Shift $=\dim(G_2)$]
The Higgs mechanism shifts $\SU(2)$ by exactly $\dim(G_2)=14$:
\begin{equation}
\boxed{\alpha_2^{-1} = \alpha_2^{0\,-1} + \dim(G_2) = 34 + 14 = 48 = 2f}
\end{equation}
where $\dim(G_2)=\lambda\Phisix = 2\times 7 = 14$ and $G_2=\Aut(\mathbb{O})$.
\end{theorem}

\begin{corollary}[Higgs Mass from $G_2$ Shift]
The Higgs quartic $\lambda_H = \Phisix/(2q^3)=7/54$ gives:
\begin{equation}
m_H = \vEW\sqrt{2\lambda_H} = 246\,\sqrt{\frac{14}{54}} = 125.3\;\mathrm{GeV}
\end{equation}
Measured: $m_H = 125.25\pm 0.17$~GeV --- deviation \textbf{0.04\%}.
\end{corollary}


\section{The Koide--Coupling Unification}
\label{sec:koide-coupling}

\begin{theorem}[The $\lambda/q$ Master Principle]
The ratio $\lambda/q = 2/3$ simultaneously governs:
\begin{enumerate}[nosep]
\item \textbf{Masses}: Koide formula
  $Q = \sum m_i/(\sum\sqrt{m_i})^2 = \lambda/q = 2/3$ (leptons, $10^{-5}$ verified);
\item \textbf{Couplings}: $\alpha_3(M_{\mathrm{Pl}})/\alpha_1(M_{\mathrm{Pl}}) = 34/51 = \lambda/q = 2/3$;
\item \textbf{Spectral data}: eigenvalue/field ratio $r/q = \lambda/q = 2/3$.
\end{enumerate}
\end{theorem}

The coupling chain factors through $17 = \Phithree+\mu = \Theta+\Phisix$:
\begin{equation}
34 \;\xrightarrow{\;\times\,f/17\;}\; 48
\;\xrightarrow{\;\times\,17/s^2\;}\; 51, \qquad
\text{product} = \frac{f}{s^2} = \frac{24}{16} = \frac{q}{\lambda} = \frac{3}{2}.
\end{equation}


\section{The Exceptional Chain as Symmetry Breaking}
\label{sec:exceptional-chain}

\begin{theorem}[E${}_8\to$\,SM via the Exceptional Series]
The five exceptional Lie algebras form the complete symmetry
breaking chain from $E_8$ to the Standard Model:
\begin{equation}
\boxed{
E_8 \xrightarrow{-115} E_7 \xrightarrow{-55} E_6
\xrightarrow{-26} F_4 \xrightarrow{-38} G_2 \xrightarrow{-14} \mathrm{SM}
}
\end{equation}
with $115+55+26+38+14 = 248 = \dim(E_8)$, and each step having a
W(3,3) formula and physical interpretation:
\end{theorem}

\begin{center}
\begin{tabular}{@{}clll@{}}
\toprule
\textbf{Step} & \textbf{Loss} & \textbf{W(3,3) formula} & \textbf{Physics} \\
\midrule
$E_8\to E_7$ & 115 & $(\mu+1)(\Phithree+\Theta)$ & gravity separation \\
$E_7\to E_6$ & 55 & $\binom{k-1}{2} = N_{\mathrm{eff}}$ & unified coupling \\
$E_6\to F_4$ & 26 & $\lambda\Phithree$ & bosonic string dim \\
$F_4\to G_2$ & 38 & $\lambda(k+\Phisix)$ & projective structure \\
$G_2\to\mathrm{SM}$ & 14 & $\lambda\Phisix = \dim(G_2)$ & EWSB \\
\bottomrule
\end{tabular}
\end{center}

The bottom three steps sum to $26+38+14 = 78 = \dim(E_6) = \lambda(v{-}1)$.


\section{The Physical Meaning of the Fine-Structure Constant}
\label{sec:alpha-meaning}

\begin{theorem}[$\alpha^{-1}$ as Geometry Plus Degrees of Freedom]
\label{thm:alpha-dof}
\begin{equation}
\boxed{\alpha^{-1} = 137 = \underbrace{\Phithree}_{13} +
\underbrace{\mu(f+\Phisix)}_{124} =
\text{(local geometry)} + \text{(total SM DOF)}}
\end{equation}
where the $124$ physical degrees of freedom are:
\begin{equation}
124 = \underbrace{27}_{\substack{\text{gauge boson}\\\text{DOF}}} +
\underbrace{1}_{\text{Higgs}} +
\underbrace{96}_{\substack{\text{fermion}\\\text{DOF}}}
= q^3 + 1 + 2q\cdot s^2.
\end{equation}
\end{theorem}

\begin{theorem}[Gauge Boson DOF $= q^3 = \dim(E_6\text{ fund.})$]
After EWSB, the physical gauge boson polarisation count is:
\begin{equation}
\underbrace{s^2}_{16\;\text{(gluons)}} +
\underbrace{q!}_{6\;\text{($W^\pm$)}} +
\underbrace{q}_{3\;\text{($Z$)}} +
\underbrace{\lambda}_{2\;\text{($\gamma$)}} = 27 = q^3.
\end{equation}
\end{theorem}

\begin{corollary}[Connection to Moonshine]
The same $124 = \mu(f+\Phisix)$ appears in the $j$-function grade-2 closure:
\begin{equation}
2099 = 244 + g\cdot 124 - (q+2).
\end{equation}
\end{corollary}
